% -----------------------------------------------------------
\section{Other Extended Scripture Studies}
% -----------------------------------------------------------
	\label{EXSS-OTHER}
	
	% ----------------------
	\subsection{Joseph Smith}
	% ----------------------
		\label{EXSS-OTHER-JS}
		
		% -----
		\subsubsection{On the meaning of the term ``Mormon''}
		% -----
			\label{EXSS-OTHER-JS-MORMON}
			
			In a \hyperref[EMS-1843-JS-CIR-20-MAY-1843]{letter to the editor} of the \textit{Times and Seasons} written in May 1843, JS comments on the meaning of the term ``Mormon,'' and says that it means ``more good.'' 
			
			Here is a discussion of that term and JS's comments, from the \href{https://onoma.lib.byu.edu/index.php/MORMON}{Book of Mormon Onomasticon} (pulled 27 June 2022), for the entry ``Mormon.'' Some of this seems pretty sloppy to me.
			
			``Any discussion of the meaning of the name \textbf{\textsc{MORMON}}, must start by dealing with a letter published in the \textit{Times and Seasons} 4 (15 May 1843): 194, that is attributed to the Prophet Joseph Smith.

			``But first a probable explanation of why the letter was written in the first place is in order. In E. D. Howe's 1834 \textit{Mormonism Unvailed}, it is claimed that `the word \textit{Mormon}, the name given to his [Joseph Smith's] book, is the English termination of the Greek word ``\textit{Mormoo},'' which we find defined in an old, obsolete Dictionary, to mean ``\textit{bug-bear, hob-goblin, raw head, and bloody bones}.''' Almost any knowledgeable reader, even in 1834, would have recognized that this definition is not only fabricated but downright silly. Closer in time to the letter in question is this passage from a local Illinois newspaper in 1841: `I will here give you the signification of the word Mormon, and also, book of Mormon, which every person that has read a dictionary of the reformed Egyptian tongue knows to be correct. \textit{Mormon}---A writer of wicked, absurd, fictitious nonsense, for evil purposes, to make sorcerers. \textit{Book of Mormon}---A book of gross, fictitious nonsense, wrote by Mormon, for Gazelom's diabolical purposes. \textit{Mormons}---Anciently in Egypt---a set of black-legs, thieves, robbers, and murderers.' This satirical attempt to define \textit{Mormon} is even more fanciful and absurd than E. D. Howe's. Such doggerel regarding \textit{Mormon} became the standard fare in the yellow journalism of the times. But no matter how outdated and fetid the nonsense, a reply seems to have been the reason for writing the letter that was published in 1843 in the \textit{Times and Seasons}.

			``And now the letter, which was printed over the name of the Prophet: `I may safely say that the word Mormon stands independent of the learning and wisdom of this generation.---Before I give a definition, however, to the word, let me say that the Bible in its widest sense, means \textit{good}; for the Savior says according to the gospel of John, ``I am the good shepherd;'' and it will not be beyond the common use of terms, to say that good is among the most important in use, and though known by various names in different languages, still its meaning is the same, and is ever in opposition to \textit{bad}. We say from the Saxon, \textit{good}; the Dane, \textit{god}; the Goth, \textit{goda}; the German, \textit{gut}; the Dutch, \textit{goed}; the Latin, \textit{bonus}; the Greek, \textit{kalos}; the \textsc{HEBREW}, \textit{tob}; and the Egyptian, \textit{mon}. Hence, with the addition of more, or the contraction, \textit{mor}, we have the word \textbf{\textsc{MORMON}}; which means, literally, \textit{more good}' (\textit{T\&S 4} [15 May 1843]: 194).

			``It is possible that the tone of the letter was at least partially meant to imitate the flippant anti-Mormon literature of the previous ten years. After all, satire is a tempting retort to satire. Although some of the letter might be an application of \textit{lex talionis} (an eye for an eye), there is a more salient crux that needs to be addressed.

			``The first issue with this statement is that it is not certain Joseph Smith is responsible for all of the content. The Prophet's journal entry for 20 May 1844 [1843] reads, `in the office heard Bro Phelps read a deffinition of the Word of Mormon -- More-Good -- corrected and sent to press.' Unfortunately, not enough information is given to determine which parts of the letter published over Joseph's signature stem from W. W. Phelps and which parts Joseph corrected. What is certain is that Joseph was not the original author, but that Joseph made changes in the text, and that he gave approval to have it published over his name. This was not the first or last time that W. W. Phelps was a ghostwriter for Joseph. B. H. Roberts, when he was compiling the \textit{Documentary History of the Church} (hereafter \textit{HC}), also `found evidence that the editor of \textit{Times and Seasons}, W. W. Phelps, rather than Joseph Smith, wrote this paragraph and that it was ``based on inaccurate premises and was offensively pedantic.''' He asked for and received permission from the First Presidency to leave the offending paragraph out of the official \textit{History of the Church} he was producing. [The footnote here is: ``Truman Madsen, \textit{Defender of the Faith: The B. H. Roberts Story} (Salt Lake City: Bookcraft, 1980), 291-292. Madsen cites no source for his information.''] In the final version of the \textit{HC}, B. H. Roberts introduced the letter with a paraphrase of Joseph''s journal entry, rather frankly writing, `Corrected and sent to the \textit{Times and Seasons} the following.' After leaving out all the words after `the learning and wisdom of this generation,' Roberts summarized the last sentence as `The word Mormon, means literally, more good.'

			``But there are other issues with this letter. When the letter states that `the Bible in its widest sense, means \textit{good},' the writer was not suggesting that the word `Bible' etymologically means \textit{good}. Rather, the writer was suggesting that the Bible is good and the reading of it promotes good. This certainly is an acceptable metaphorical meaning of `Bible' that no Christian in the 19th century would deny.

			``This metaphorical meaning leads to an examination of the phrase, `The word Mormon, means literally, more good.' Just as today, the word \textit{literally} was used in 19th century English in the sense of `actually,' `really.'[7] In other words, the word `literally' can be understood in the letter to mean `actually' and not necessarily `word-for-word.' Thus, if the Bible, which is true as far as it was transmitted correctly, means `good,' then the Book of Mormon, which was transmitted correctly, must actually mean `more good.' The conclusion can be drawn that `more good' is not a translation of the word `Mormon,' but a metaphorical interpretation in the sense that while the Bible means `good, `Mormon means `more [of the] good,' or possibly `better translated than the Bible.'

			``What then is the philological explanation of \textbf{\textsc{MORMON}}? Notwithstanding the warning in the letter attributed to Joseph Smith `that the word Mormon stands independent of the learning and wisdom of this generation,' many attempts have been made to provide a sound etymology for \textbf{\textsc{MORMON}} based on secular knowledge of ancient Near Eastern languages. The results at best show promise. The following discussion reviews some of the suggestions that have been made.

			``The first point to be made is that the name, being mentioned first as a GN [geographical name] (Mosiah 18:4) and then later as a PN [personal name] (3 Nephi 5:12, in which the PN is explicitly derived from the GN), might derive from a descriptive that would be appropriate for both a place and a person. It might be suggested, based on Mosiah 18:4, `a place which was called Mormon, having received its name from the king ... having been infested... by wild beasts,' that \textbf{\textsc{MORMON}} could have something to do with `wild beasts.' On the other hand, based on Mosiah 18:5, `there was in Mormon a fountain of pure water,' or based on Mosiah 18:30, `the place of Mormon, the waters of Mormon, the forest of Mormon, how beautiful are they,' \textbf{\textsc{MORMON}} might derive from `fountain/spring' or `pure water' or `beautiful/beauty.'

			``Janne Sjodahl derived the second half of the word MORMON from ancient Egyptian \textit{'Imn} `Amun,' and Hugh Nibley similarly suggested that MORMON possibly derived from Egyptian PN \textit{Mry-`Imn}, `Beloved of ``Amon' -- applying Egyptian \textit{mr}- `want, wish, desire, love; beloved,' along with the name of the Egyptian god \textit{'Amun}. Meir Lubetski called our attention to pre-exilic Hebrew use of this Egyptian name \textit{Mry-`Imn} on Hebrew seals as \textit{my'mn}. This makes sense since the -\textit{r}- in \textit{mr(y)} is dropped in Egyptian scribal Akkadian \textit{ma-a-i}, as well as in Coptic \textit{me}, \textit{mei}, \textit{mai}, \textit{m\={e}i} `love.' Matt Bowen demonstrates how this word is used for several puns (Mosiah 18:8--11,21,28,30, Alma 5:3--6,26, Mormon 3:12, Moroni 7:47--48, 8:16--17), and he applies the same understanding of the etymology to the name MARY in the Book of Mormon.

			``Similarly, Robert F. Smith and Benjamin Urrutia have each called attention to \textsc{EGYPTIAN} \textit{mr mn}, `truly beloved,' or `love is established' (\textit{BU},), or `strong/firm love' or `love remains steadfast/firm' (RFS). The translation `love is established forever' brings to memory the words of Paul, `charity never faileth' (1 Corinthians 13:8) (\textit{BU}). Interestingly, it is \textbf{\textsc{MORMON}} who uses the same words in a letter written to his son MORONI, adding, `But charity is the pure love of CHRIST, and it endureth forever' (Moroni 7:46--47) (JAT).

			``As Carl Hugh Jones pointed out in 1970, Frederick G. Williams left two possible hieratic characters from the Book of Mormon plates which he claimed represent the phrase `Book of Mormon,' and Egyptologist John Gee has said that it looks authentic. Why? Because the character for `book' in that phrase is written nearly identical to a known early Demotic Egyptian form of \textit{m\_{d}3t} `book.' That the syntax appears to be \textit{Mrmn md3t} (reading from right to left) is no problem since names or titles frequently come first in Egyptian, as in \textit{m\_{d}3t n\_{t}r} `Holy Writ,' where \textit{n\_{t}r} is written first.

			``On a limestone stele of the 19th to 21st \textsc{EGYPTIAN} dynasty in the Museum of Gizeh the name \textit{mrmnu} appears, accompanied by the title `door keeper.' In a yet to be superseded article, W. Spiegelberg, `Zu den semitischen Eigennamen in ägyptischer Umschrift aus der Zeit des ``neuen Reiches'' (um 1500--1000),' \textit{Zeitschrift für Assyriologie} 13 (1898):51, treats the name as Semitic in \textsc{EGYPTIAN} transcription although he is not certain that it is Semitic, and he does not provide a meaning. He transcribes it into \textsc{HEBREW} characters with \textit{mr/lmn(w)}. Spiegelberg's description of the stele unfortunately does not permit its current identification. Despite various difficulties, such as dating to at least 600 years before \textsc{LEHI} and not having an etymology, this name, \textit{mrmn}, on an \textsc{EGYPTIAN} inscription seems like a direct hit, as Hugh Nibley pointed out years ago. Unfortunately, no immediate etymology suggests itself.

			``Nevertheless, Nibley has also pointed out that \textit{mrm}, besides appearing in the \textsc{EGYPTIAN} PN, is attested in \textsc{HEBREW} and Arabic, and means `desirable' or `good.' In this case, \textbf{\textsc{MORMON}} would consist of the root \textit{mrm} plus the common Semitic ending \textit{-\={o}n} (often used on GNs and PNs, such as Kidron and \textsc{GIDEON}). For possible, \textsc{HEBREW} examples, see the biblical PN \textit{mirm\^{a}h} in 1 Chronicles 8:10 (\textit{HALOT} does not offer an etymology; the Septuagint transcribes it as \gr{μαρμα}), and the PN \textit{mer\={e}m\^{o}t} (also of questionable etymology; Septuagint \gr{μεραμωθ}), the name of a priest in Ezra 10:36 (=Nehemiah 10:5), are possibilities (JH). The name also appears as \textit{mrmwt} on a 6th century BC ostracon from Arad. Note also the PN at Ugarit, \textit{ma-ri-ma-na} (JH), but the language origin of the name is unknown.

			``Less likely is \textsc{EGYPTIAN} \textit{mr} (> Nubian and Coptic \textit{mur}, \textit{mor}), `bind, girth' (RFS). On the element \textit{mr} see Nibley, \textit{SC}, 194, n. 107.''